% ----------------------------------------------------
%
%            Source Verse, Essential Verse
%        and desugaring from the former to the latter
%
% ----------------------------------------------------

%% ODER: format ==         = "\mathrel{==}"
%% ODER: format /=         = "\neq "
%
%
\makeatletter
\@ifundefined{lhs2tex.lhs2tex.sty.read}%
  {\@namedef{lhs2tex.lhs2tex.sty.read}{}%
   \newcommand\SkipToFmtEnd{}%
   \newcommand\EndFmtInput{}%
   \long\def\SkipToFmtEnd#1\EndFmtInput{}%
  }\SkipToFmtEnd

\newcommand\ReadOnlyOnce[1]{\@ifundefined{#1}{\@namedef{#1}{}}\SkipToFmtEnd}
\usepackage{amstext}
\usepackage{amssymb}
\usepackage{stmaryrd}
\DeclareFontFamily{OT1}{cmtex}{}
\DeclareFontShape{OT1}{cmtex}{m}{n}
  {<5><6><7><8>cmtex8
   <9>cmtex9
   <10><10.95><12><14.4><17.28><20.74><24.88>cmtex10}{}
\DeclareFontShape{OT1}{cmtex}{m}{it}
  {<-> ssub * cmtt/m/it}{}
\newcommand{\texfamily}{\fontfamily{cmtex}\selectfont}
\DeclareFontShape{OT1}{cmtt}{bx}{n}
  {<5><6><7><8>cmtt8
   <9>cmbtt9
   <10><10.95><12><14.4><17.28><20.74><24.88>cmbtt10}{}
\DeclareFontShape{OT1}{cmtex}{bx}{n}
  {<-> ssub * cmtt/bx/n}{}
\newcommand{\tex}[1]{\text{\texfamily#1}}	% NEU

\newcommand{\Sp}{\hskip.33334em\relax}


\newcommand{\Conid}[1]{\mathit{#1}}
\newcommand{\Varid}[1]{\mathit{#1}}
\newcommand{\anonymous}{\kern0.06em \vbox{\hrule\@width.5em}}
\newcommand{\plus}{\mathbin{+\!\!\!+}}
\newcommand{\bind}{\mathbin{>\!\!\!>\mkern-6.7mu=}}
\newcommand{\rbind}{\mathbin{=\mkern-6.7mu<\!\!\!<}}% suggested by Neil Mitchell
\newcommand{\sequ}{\mathbin{>\!\!\!>}}
\renewcommand{\leq}{\leqslant}
\renewcommand{\geq}{\geqslant}
\usepackage{polytable}

%mathindent has to be defined
\@ifundefined{mathindent}%
  {\newdimen\mathindent\mathindent\leftmargini}%
  {}%

\def\resethooks{%
  \global\let\SaveRestoreHook\empty
  \global\let\ColumnHook\empty}
\newcommand*{\savecolumns}[1][default]%
  {\g@addto@macro\SaveRestoreHook{\savecolumns[#1]}}
\newcommand*{\restorecolumns}[1][default]%
  {\g@addto@macro\SaveRestoreHook{\restorecolumns[#1]}}
\newcommand*{\aligncolumn}[2]%
  {\g@addto@macro\ColumnHook{\column{#1}{#2}}}

\resethooks

\newcommand{\onelinecommentchars}{\quad-{}- }
\newcommand{\commentbeginchars}{\enskip\{-}
\newcommand{\commentendchars}{-\}\enskip}

\newcommand{\visiblecomments}{%
  \let\onelinecomment=\onelinecommentchars
  \let\commentbegin=\commentbeginchars
  \let\commentend=\commentendchars}

\newcommand{\invisiblecomments}{%
  \let\onelinecomment=\empty
  \let\commentbegin=\empty
  \let\commentend=\empty}

\visiblecomments

\newlength{\blanklineskip}
\setlength{\blanklineskip}{0.66084ex}

\newcommand{\hsindent}[1]{\quad}% default is fixed indentation
\let\hspre\empty
\let\hspost\empty
\newcommand{\NB}{\textbf{NB}}
\newcommand{\Todo}[1]{$\langle$\textbf{To do:}~#1$\rangle$}

\EndFmtInput
\makeatother
%
%
%
%
%
%
% This package provides two environments suitable to take the place
% of hscode, called "plainhscode" and "arrayhscode". 
%
% The plain environment surrounds each code block by vertical space,
% and it uses \abovedisplayskip and \belowdisplayskip to get spacing
% similar to formulas. Note that if these dimensions are changed,
% the spacing around displayed math formulas changes as well.
% All code is indented using \leftskip.
%
% Changed 19.08.2004 to reflect changes in colorcode. Should work with
% CodeGroup.sty.
%
\ReadOnlyOnce{polycode.fmt}%
\makeatletter

\newcommand{\hsnewpar}[1]%
  {{\parskip=0pt\parindent=0pt\par\vskip #1\noindent}}

% can be used, for instance, to redefine the code size, by setting the
% command to \small or something alike
\newcommand{\hscodestyle}{}

% The command \sethscode can be used to switch the code formatting
% behaviour by mapping the hscode environment in the subst directive
% to a new LaTeX environment.

\newcommand{\sethscode}[1]%
  {\expandafter\let\expandafter\hscode\csname #1\endcsname
   \expandafter\let\expandafter\endhscode\csname end#1\endcsname}

% "compatibility" mode restores the non-polycode.fmt layout.

\newenvironment{compathscode}%
  {\par\noindent
   \advance\leftskip\mathindent
   \hscodestyle
   \let\\=\@normalcr
   \let\hspre\(\let\hspost\)%
   \pboxed}%
  {\endpboxed\)%
   \par\noindent
   \ignorespacesafterend}

\newcommand{\compaths}{\sethscode{compathscode}}

% "plain" mode is the proposed default.
% It should now work with \centering.
% This required some changes. The old version
% is still available for reference as oldplainhscode.

\newenvironment{plainhscode}%
  {\hsnewpar\abovedisplayskip
   \advance\leftskip\mathindent
   \hscodestyle
   \let\hspre\(\let\hspost\)%
   \pboxed}%
  {\endpboxed%
   \hsnewpar\belowdisplayskip
   \ignorespacesafterend}

\newenvironment{oldplainhscode}%
  {\hsnewpar\abovedisplayskip
   \advance\leftskip\mathindent
   \hscodestyle
   \let\\=\@normalcr
   \(\pboxed}%
  {\endpboxed\)%
   \hsnewpar\belowdisplayskip
   \ignorespacesafterend}

% Here, we make plainhscode the default environment.

\newcommand{\plainhs}{\sethscode{plainhscode}}
\newcommand{\oldplainhs}{\sethscode{oldplainhscode}}
\plainhs

% The arrayhscode is like plain, but makes use of polytable's
% parray environment which disallows page breaks in code blocks.

\newenvironment{arrayhscode}%
  {\hsnewpar\abovedisplayskip
   \advance\leftskip\mathindent
   \hscodestyle
   \let\\=\@normalcr
   \(\parray}%
  {\endparray\)%
   \hsnewpar\belowdisplayskip
   \ignorespacesafterend}

\newcommand{\arrayhs}{\sethscode{arrayhscode}}

% The mathhscode environment also makes use of polytable's parray 
% environment. It is supposed to be used only inside math mode 
% (I used it to typeset the type rules in my thesis).

\newenvironment{mathhscode}%
  {\parray}{\endparray}

\newcommand{\mathhs}{\sethscode{mathhscode}}

% texths is similar to mathhs, but works in text mode.

\newenvironment{texthscode}%
  {\(\parray}{\endparray\)}

\newcommand{\texths}{\sethscode{texthscode}}

% The framed environment places code in a framed box.

\def\codeframewidth{\arrayrulewidth}
\RequirePackage{calc}

\newenvironment{framedhscode}%
  {\parskip=\abovedisplayskip\par\noindent
   \hscodestyle
   \arrayrulewidth=\codeframewidth
   \tabular{@{}|p{\linewidth-2\arraycolsep-2\arrayrulewidth-2pt}|@{}}%
   \hline\framedhslinecorrect\\{-1.5ex}%
   \let\endoflinesave=\\
   \let\\=\@normalcr
   \(\pboxed}%
  {\endpboxed\)%
   \framedhslinecorrect\endoflinesave{.5ex}\hline
   \endtabular
   \parskip=\belowdisplayskip\par\noindent
   \ignorespacesafterend}

\newcommand{\framedhslinecorrect}[2]%
  {#1[#2]}

\newcommand{\framedhs}{\sethscode{framedhscode}}

% The inlinehscode environment is an experimental environment
% that can be used to typeset displayed code inline.

\newenvironment{inlinehscode}%
  {\(\def\column##1##2{}%
   \let\>\undefined\let\<\undefined\let\\\undefined
   \newcommand\>[1][]{}\newcommand\<[1][]{}\newcommand\\[1][]{}%
   \def\fromto##1##2##3{##3}%
   \def\nextline{}}{\) }%

\newcommand{\inlinehs}{\sethscode{inlinehscode}}

% The joincode environment is a separate environment that
% can be used to surround and thereby connect multiple code
% blocks.

\newenvironment{joincode}%
  {\let\orighscode=\hscode
   \let\origendhscode=\endhscode
   \def\endhscode{\def\hscode{\endgroup\def\@currenvir{hscode}\\}\begingroup}
   %\let\SaveRestoreHook=\empty
   %\let\ColumnHook=\empty
   %\let\resethooks=\empty
   \orighscode\def\hscode{\endgroup\def\@currenvir{hscode}}}%
  {\origendhscode
   \global\let\hscode=\orighscode
   \global\let\endhscode=\origendhscode}%

\makeatother
\EndFmtInput
%
%
%
% First, let's redefine the forall, and the dot.
%
%
% This is made in such a way that after a forall, the next
% dot will be printed as a period, otherwise the formatting
% of `comp_` is used. By redefining `comp_`, as suitable
% composition operator can be chosen. Similarly, period_
% is used for the period.
%
\ReadOnlyOnce{forall.fmt}%
\makeatletter

% The HaskellResetHook is a list to which things can
% be added that reset the Haskell state to the beginning.
% This is to recover from states where the hacked intelligence
% is not sufficient.

\let\HaskellResetHook\empty
\newcommand*{\AtHaskellReset}[1]{%
  \g@addto@macro\HaskellResetHook{#1}}
\newcommand*{\HaskellReset}{\HaskellResetHook}

\global\let\hsforallread\empty

\newcommand\hsforall{\global\let\hsdot=\hsperiodonce}
\newcommand*\hsperiodonce[2]{#2\global\let\hsdot=\hscompose}
\newcommand*\hscompose[2]{#1}

\AtHaskellReset{\global\let\hsdot=\hscompose}

% In the beginning, we should reset Haskell once.
\HaskellReset

\makeatother
\EndFmtInput
% ----------------------
%% Stuff for TimCore


%% 'g' inferences





%% Intersection of effects

%% --------- Effects ------------


%% ----------------------


%% --------------------------
%%        Desugaring rules
%% --------------------------



% ----------------------------------------------------
% New desugaring, just does the wrapping part:
% ----------------------------------------------------
%   Input for wrapping


%%format Hide (e) = "{\mathcal H}\llbracket{}" e "\rrbracket{}"

%% 'si' is short for "solve or infer" or "unification variable or not" resp

%% 'md' is short for "mode" or "unification variable or not" resp

%% flipsi is a no-op for now


%% EXX is just EX, but with no argument




%% ----------- Applications --------------
%% %format applyx a b = a "@" b
%% %format applyx a b = a "\," b

%% Function application with no arguments

%% Function application with 1 argument

%% Function application with 1 argument, no parens

%% Function application with 1 argument, without parenthesis

%% Function application with many arguments

%% ----------- End of Applications --------------





%% Sequences e1; ...; en; v

%% xs etc stole syntax used in examples.

%% Old version: | for BNF, [] for choice
%%%format choice = "[ \! ]"
%%%format `choice` = "\mathop{[ \! ]}"
%% 1mm is about 0.3em with the current font

%% New version: tall | for BNF, fat | for choice
%%format vbar = "\mathop{\bigm|}"
%%format choice = "\boldsymbol{|}"
%%format `choice` = "\mathop{\boldsymbol{|}}"




%% %format ok = "\textbf{ok}"

%% %format `eqsemi` = "\hspace{-0.5mm}\fatsemi{}"
%% %format eqsemiKW = "\fatsemi"

%% Arrays and tuples

%% %format vmap (x) = "\llangle" x "\rrangle"






%%  OLD SYNTAX   %format split (x) (y) (z) = "\textbf{split}(" x ")\{" y "," z "\}"
%% %format defKW = "\textbf{def}"
% only used in comments, but still prettier this way:
%%%%format map = "\textbf{map}"

%% -------Functions ---------------


%% functionOC has open/closed subscript


%% -------End of f unctions ---------------





%% Unification
%%%format == = "\!=\!"


% format := = "\mapsto"
% %format squiggt = "\leadsto"
% %format squiggt = "\!-\!\!>\!\!\!"



































%% We combined S and C into one


%% ----------- Metatheory --------------


%% ----------------------------
%%         Verification
%% ----------------------------

%% assert = succeed, abbreviated to suc



%% Effects checking, like succeeds{e}





%% <fx> brackets

\subsection{Source Verse}

\begin{figure}
$$
\begin{array}{lrcll}
\textbf{Terms} & \ensuremath{\Varid{t}} & ::= & \ensuremath{\Varid{x}}               & \text{Variable} \\
&      & \vb & \ensuremath{\anonymous }               & \text{Wildcard} \\
&      & \vb & \ensuremath{\Varid{k}}               & \text{Constant} \\
&      & \vb & \ensuremath{\Varid{t}_{\mathrm{1}};\,\Varid{t}_{\mathrm{2}}}       & \text{Sequencing} \\
&      & \vb & \ensuremath{\Varid{t}_{\mathrm{1}} [\mskip1.5mu \Varid{t}_{\mathrm{2}}\mskip1.5mu]}       & \text{Function application} \\
&      & \vb & \ensuremath{\Varid{t}_{\mathrm{1}} (\Varid{t}_{\mathrm{2}})}       & \text{Non-failing function application} \\
&      & \vb & \ensuremath{\exists\Varid{x}.\,\Varid{t}}         & \text{Existential} \\
&      & \vb & \ensuremath{\Varid{t}_{\mathrm{1}}\hspace{0.19em}\mathop{\rule[-0.1em]{0.15em}{0.75em}}\hspace{0.2em}\Varid{t}_{\mathrm{2}}} & \text{Choice} \\
&      & \vb & \ensuremath{\Varid{t}_{\mathrm{1}}\mathrel{=}\Varid{t}_{\mathrm{2}}}       & \text{Unification} \\
&      & \vb & \ensuremath{\textbf{if}\;(\Varid{t}_{\mathrm{1}})\;\textbf{then}\;\Varid{t}_{\mathrm{2}}\;\textbf{else}\;\Varid{t}_{\mathrm{3}}} & \text{Conditional expression} \\
&      & \vb & \ensuremath{\textbf{for}\;(\Varid{t}_{\mathrm{1}})\;\textbf{do}\;\Varid{t}_{\mathrm{2}}}     & \text{Finite loop} \\
&      & \vb & \ensuremath{{\bf chk}\langle\omega\rangle\{\Varid{t}\}}     & \text{Effects check} \\
&      & \vb & \ensuremath{\textbf{array}\{\Varid{t}_{\mathrm{1}};\,\mydots;\,\Varid{t}_{\Varid{n}}\}} & \text{Array}, n \geq 0 \\
&      & \vb & \ensuremath{\Varid{t}_{\mathrm{1}},\mydots,\Varid{t}_{\Varid{n}}}  & \text{Array}, n \geq 2 \\
&      & \vb & \ensuremath{()}              & \text{Empty array} \\
&      & \vb & \ensuremath{\Varid{t}_{\mathrm{1}}\;\!=>\!\!\;\Varid{t}_{\mathrm{2}}}    & \text{Lambda expression} \\
&      & \vb & \ensuremath{\Varid{p}\mathrel{\coloneqq}\Varid{t}}          & \text{General definition} \\
&      & \vb & \ensuremath{\Varid{p}\mathbin{:}\Varid{t}}           & \text{Type ascription} \\
&      & \vb & \ensuremath{\Varid{op}_{\Varid{b}}\;\Varid{t}}          & \text{Prefix operator} \\
&      & \vb & \ensuremath{\Varid{t}\;\Varid{op}_{\Varid{a}}}          & \text{Postfix operator} \\
&      & \vb & \ensuremath{\Varid{t}_{\mathrm{1}}\;\Varid{op}_{\Varid{i}}\;\Varid{t}_{\mathrm{2}}}    & \text{Infix operator} \\
&      & \vb & \ensuremath{\mathbin{:}\Varid{t}}             & \text{Range} \\
&      & \vb & \ensuremath{\Varid{t}_{\mathrm{1}}\;\mathbf{where}\;\Varid{t}_{\mathrm{2}}}   & \text{Reverse sequencing} \\
&      & \vb & \ensuremath{\textbf{fun}_{\Varid{q}}(\Varid{t})\langle\omega\rangle\{\Varid{t}\}}    & \text{Function} \\
&      & \vb & \ensuremath{\textbf{if}\;(\Varid{t})\;\textbf{then}\;\Varid{t}}    & \text{Conditional expression} \\
&      & \vb & \ensuremath{\textbf{if}\;(\Varid{t})\;\textbf{else}\;\Varid{t}}    & \text{Conditional expression} \\
&      & \vb & \ensuremath{\textbf{if}\;\{\mskip1.5mu \Varid{t}\mskip1.5mu\}}           & \text{Conditional expression} \\
&      & \vb & \ensuremath{\textbf{for}\;\{\mskip1.5mu \Varid{t}\mskip1.5mu\}}          & \text{Finite loop} \\
&      & \vb & \ensuremath{\mathbf{let}\;(\Varid{t})\;\textbf{in}\;\Varid{t}}     & \text{Local binding} \\
&      & \vb & \ensuremath{\textbf{block}\;\Varid{t}}         & \text{Scope delimiter} \\
&      & \vb & \ensuremath{\textbf{case}\;(\Varid{t})\;\textbf{of}\;\{\mskip1.5mu \Varid{t}_{\mathrm{1}};\,\mydots\;\Varid{t}_{\Varid{n}}\mskip1.5mu\}}    & \text{Case expression} \\
&      & \vb & \ensuremath{\textbf{case}\;\{\mskip1.5mu \Varid{t}_{\mathrm{1}};\,\mydots\;\Varid{t}_{\Varid{n}}\mskip1.5mu\}} & \text{Case expression} \\
&      & \vb & \ensuremath{\textbf{type} \{\mskip1.5mu \Varid{t}\mskip1.5mu\}}         & \text{Type definition} \\
&      & \vb & \ensuremath{\textbf{option} \{\mskip1.5mu \mskip1.5mu\}}        & \text{Create an option} \\
&      & \vb & \ensuremath{\textbf{option} \{\mskip1.5mu \Varid{t}\mskip1.5mu\}}       & \text{Create an option} \\
&      & \vb & \ensuremath{\textbf{truth}\{\Varid{t}\}}       & \text{Singleton map} \\
&      & \vb & \ensuremath{\textbf{map}\{\Varid{t}_{\mathrm{1}};\,\mydots;\,\Varid{t}_{\Varid{n}}\}} & \text{Map}, n \geq 0 \\
\\
\text{\bf Operators} &
 \ensuremath{\Varid{op}_{\Varid{b}}} & ::= &  ! \vb - \vb + \vb \ensuremath{\hat{\hspace{0.8em}}} \vb ? \vb \ensuremath{[]} \vb ..  & \text{Prefix}\\
& \ensuremath{\Varid{op}_{\Varid{a}}} & ::= &  ? \vb \ensuremath{\hat{\hspace{0.8em}}} & \text{Postfix}\\
& \ensuremath{\Varid{op}_{\Varid{i}}} & ::= &  + \vb - \vb * \vb / \vb \ensuremath{\textbf{and}} \vb \ensuremath{\textbf{or}} \vb : \vb .. & \text{Infix} \\
&        & \vb & < \vb <= \vb > \vb >= \vb <> \vb \ensuremath{-\!\!>} \\
\end{array}
$$
\caption{Source Verse expressions, part 1} \label{fig:source-verse-expr}
\end{figure}

\begin{figure}
$$
\begin{array}{lrcll}
  \text{\bf Patterns} \asktim & \ensuremath{\Varid{p}} & ::= & \ensuremath{\anonymous }   & \text{Wildcard} \\
  & & \vb & \ensuremath{ xs }                         & \text{Variable} \\
  & & \vb & \ensuremath{\Varid{p}_{\mathrm{1}}\!\rightarrow\!\Varid{p}_{\mathrm{2}}}               & \text{Squiggle pattern} \\
  & & \vb & \ensuremath{\Varid{p}(\Varid{t})\;\langle\omega\rangle}       & \text{Function with effects} \\
  & & \vb & \ensuremath{\Varid{p}(\Varid{t})}              & \text{Short for \ensuremath{\Varid{p}(\Varid{t})\;\langle\textbf{succeeds}\rangle} } \\
  & & \vb & \ensuremath{\Varid{p}\mathbin{:}\Varid{t}\;\langle\omega\rangle}              & \text{Type ascription}\\
  & & \vb & \ensuremath{\Varid{p}\mathbin{:}\Varid{t}}                       & \text{Short for \ensuremath{\Varid{p}\mathbin{:}\Varid{t}\;\langle\textbf{decides}\rangle}}\asktim\\
  & & \vb & \ensuremath{\mathbin{:}\Varid{t}} & \text{Short for \ensuremath{(\anonymous \mathbin{:}\Varid{t})}} \\
  & & \vb & \ensuremath{\Varid{p}_{\mathrm{1}},\mydots,\Varid{p}_{\Varid{n}}}  & \text{Tuple} \\
  & & \vb & \ensuremath{\mathinner{\ldotp\ldotp}\Varid{p}} & \text{Array splice} \\
& \ensuremath{ xs } & ::= & x \vb  \ensuremath{ xs \;\!\textbf{\&}\!\;\Varid{x}}
\\[2mm]
\text{\bf Effects} &
\ensuremath{\omega} & ::= & \ensuremath{\epsilon} \vb \ensuremath{\omega_{\mathrm{1}},\omega_{\mathrm{2}}} \\
&     & \vb & \ensuremath{\textbf{fails}}  \vb \ensuremath{\textbf{succeeds}} \vb \ensuremath{\textbf{decides}} \\
&     & \vb & \ldots \\

& \ensuremath{\omega_{\mathrm{1}}\cap\omega_{\mathrm{2}}} & = & \multicolumn{2}{l}{\text{Effect intersection; e.g \ensuremath{\textbf{succeeds}\cap\textbf{decides}\mathrel{=}\textbf{succeeds}}}}

\end{array}
$$
\caption{Source Verse expressions, part 2} \label{fig:source-verse-expr2}
\end{figure}

% -------------------------------------------------------------
%               Essential Verse
% -------------------------------------------------------------

\subsection{Essential Verse}

\begin{figure}
$$
\begin{array}{lrcll}
\text{Variables} &  \ensuremath{\Varid{x},\Varid{y},\Varid{f}} &     \\

\text{Constants} & \ensuremath{\Varid{k}} & ::= & \textit{integer constants} \\
                 &     & \vb & \textit{string constants} \\
\\
\text{Terms} &  \ensuremath{\Varid{t}} & ::= & \ensuremath{\Varid{x}}                        & \text{Variable} \\
&      & \vb & \ensuremath{\anonymous }                        & \text{Underscore} \\
&      & \vb & \ensuremath{\Varid{k}}                        & \text{Constant} \\
&      & \vb & \ensuremath{\Varid{o}}                        & \text{Primops} \\
&      & \vb & \ensuremath{\Varid{t}_{\mathrm{1}} [\mskip1.5mu \Varid{t}_{\mathrm{2}}\mskip1.5mu]}                 & \text{Function application} \\
&      & \vb & \ensuremath{\Varid{t}_{\mathrm{1}};\,\Varid{t}_{\mathrm{2}}}                 & \text{Sequencing} \\
&      & \vb & \ensuremath{\Varid{t}_{\mathrm{1}}\;\mathbf{where}\;\Varid{t}_{\mathrm{2}}}             & \text{Reverse sequencing} \\
&      & \vb & \ensuremath{\Varid{t}_{\mathrm{1}}\mathrel{=}\Varid{t}_{\mathrm{2}}}                 & \text{Unification} \\
&      & \vb & \ensuremath{\Varid{x}\mathrel{\coloneqq}\Varid{t}}                    & \text{Definition (binds \ensuremath{\Varid{x}} in \ensuremath{\Varid{t}}} \\
&      & \vb & \ensuremath{\Varid{x}\!\rightarrow\!\Varid{t}}                 & \text{Capture input (does not bind \ensuremath{\Varid{x}})} \\
&      & \vb & \ensuremath{\Varid{t}_{\mathrm{1}}\hspace{0.19em}\mathop{\rule[-0.1em]{0.15em}{0.75em}}\hspace{0.2em}\Varid{t}_{\mathrm{2}}}          & \text{Choice} \\
&      & \vb & \ensuremath{\textbf{if}\;(\Varid{t}_{\mathrm{1}})\;\textbf{then}\;\Varid{t}_{\mathrm{2}}\;\textbf{else}\;\Varid{t}_{\mathrm{3}}} & \text{Conditional expression} \\
&      & \vb & \ensuremath{\textbf{for}(\Varid{t}_{\mathrm{1}})\,\textbf{do}\,\Varid{t}_{\mathrm{2}}}             & \text{Finite loop} \\
&      & \vb & \ensuremath{\textbf{array}\{\Varid{t}_{\mathrm{1}},\mydots,\Varid{t}_{\Varid{n}}\}}   & \text{Array}, n \geq 0 \\
&      & \vb & \ensuremath{\textbf{truth}\{\Varid{t}\}}                  & \text{Truth value} \\
&      & \vb & \ensuremath{\textbf{fun}_{\Varid{q}}(\Varid{t}_{\mathrm{1}})\langle\omega\rangle\{\Varid{t}_{\mathrm{2}}\}}  & \text{Function} \\
&      & \vb & \ensuremath{\mathbin{:}\Varid{t}}                      & \text{Range} \\
&      & \vb & \ensuremath{\Varid{t}_{\mathrm{1}}\vartriangleright\hspace{-1.3mm} \langle\omega\rangle\,\Varid{t}_{\mathrm{2}}}       & \text{Opacity}\\
&      & \vb & \ensuremath{{\bf chk}\langle\omega\rangle\{\Varid{t}\}}                & \text{Effects check} \\
\\
\text{Primops} & o  & ::= & \ensuremath{\textbf{gt}} \vb \ensuremath{\textbf{add}} \vb \ensuremath{\textbf{int}} \\
\text{Aperture} &  \ensuremath{\Varid{q}} & ::=  & \ensuremath{\textbf{o}} \vb \ensuremath{\textbf{c}} & \text{Open vs closed functions}
\end{array}
$$
\\
Variables can be alphanumeric, but also \ensuremath{\textbf{prefix}\textquote \Varid{op}\textquote}, \ensuremath{\textbf{postfix}\textquote \Varid{op}\textquote}, and \ensuremath{\textbf{operator}\textquote \Varid{op}\textquote}.
\caption{Essential Verse} \label{fig:small-source-verse}
\end{figure}


% -------------------------------------------------------------
%       Desugaring  Source Verse --> Essential Verse
% -------------------------------------------------------------

\subsection{Desugaring Source Verse into Essential Verse}


\begin{figure}
$$
\begin{array}{@{}l}
\begin{array}{lrcll}
\multicolumn{4}{l}{\textbf{Abbreviations for \ensuremath{\Varid{p}\mathrel{\coloneqq}\Varid{e}}}\; \asktim} \\
\rulename{dscol1} &  \ensuremath{\Varid{p}\mathbin{:}\Varid{t}\mathrel{=}\Varid{e}} & \rightarrow & \ensuremath{\Varid{p}\mathbin{:}\Varid{t}\mathrel{\coloneqq}\Varid{e}} \\
\rulename{dscol2} &  \ensuremath{\mathbin{:}\Varid{t}\mathrel{=}\Varid{e}} & \rightarrow & \ensuremath{\mathbin{:}\Varid{t}\mathrel{\coloneqq}\Varid{e}} \\
\rulename{dscol3} &  \ensuremath{\Varid{p}\mathbin{:}\Varid{t}} & \rightarrow & \ensuremath{\Varid{p}\mathrel{\coloneqq}\mathbin{:}\Varid{t}} \\[2mm]

\multicolumn{4}{l}{\textbf{Rewriting \ensuremath{\Varid{p}\mathrel{\coloneqq}\Varid{t}}}} \\
\rulename{dswild1} & \ensuremath{\anonymous \mathrel{\coloneqq}\Varid{t}} & \rightarrow & \ensuremath{\Varid{t}}  \\
\rulename{dswild2} &  \ensuremath{\mathbin{:}\Varid{t}_{\mathrm{1}}\mathrel{\coloneqq}\Varid{t}_{\mathrm{2}}} & \rightarrow & \ensuremath{(\anonymous \mathbin{:}\Varid{t}_{\mathrm{1}})\mathrel{\coloneqq}\Varid{t}_{\mathrm{2}}} \\
\rulename{dsfun1} &  \ensuremath{\Varid{p}(\Varid{t}_{\mathrm{1}})\;\langle\Varid{q},\omega\rangle\mathrel{\coloneqq}\Varid{t}_{\mathrm{2}}}
    & \rightarrow & \ensuremath{\Varid{p}\mathrel{\coloneqq}\textbf{fun}_{\Varid{q}}(\Varid{t}_{\mathrm{1}})\langle\omega\rangle\{\Varid{t}_{\mathrm{2}}\}} \\
\rulename{dsfun2} &  \ensuremath{\Varid{p}(\Varid{t}_{\mathrm{1}})\mathrel{\coloneqq}\Varid{t}_{\mathrm{2}}}
  & \rightarrow & \ensuremath{\Varid{p}(\Varid{t}_{\mathrm{1}})\;\langle\textbf{open},\textbf{succeeds}\rangle\mathrel{\coloneqq}\Varid{t}_{\mathrm{2}}} \\
\rulename{dsty1} & \ensuremath{\Varid{p}\mathbin{:}\Varid{t}_{\mathrm{1}}\;\langle\omega\rangle\mathrel{\coloneqq}\Varid{t}_{\mathrm{2}}} & \rightarrow & \ensuremath{\Varid{p}\mathrel{\coloneqq}\Varid{t}_{\mathrm{2}}\vartriangleright\hspace{-1.3mm} \langle\omega\rangle\,\Varid{t}_{\mathrm{1}}}  \\
\rulename{dsty2}\asktim & \ensuremath{\Varid{p}\mathbin{:}\Varid{t}_{\mathrm{1}}\mathrel{\coloneqq}\Varid{t}_{\mathrm{2}}} & \rightarrow & \ensuremath{\Varid{p}\mathbin{:}\Varid{t}_{\mathrm{1}}\;\langle\textbf{decides}\rangle\mathrel{\coloneqq}\Varid{t}_{\mathrm{2}}}  \\
\rulename{dssquig} &  \ensuremath{(\Varid{p}_{\mathrm{1}}\!\rightarrow\!\Varid{p}_{\mathrm{2}})\mathrel{\coloneqq}\Varid{t}} & \rightarrow
    & \ensuremath{\Varid{p}_{\mathrm{1}}\mathrel{\coloneqq}(\Varid{x}\mathbin{:}\textbf{any});\,\Varid{p}_{\mathrm{2}}\mathrel{\coloneqq}(\Varid{x}\!\rightarrow\!\Varid{t})} & \text{\ensuremath{\Varid{x}} fresh} \\
% \begin{array}[t]{@@{}l@@{}}
%    |p_1 := (x_1:any); p_2 := (x_2:any);| \\
%    |(x_1 squiggt x_2) := t|
%  \end{array}  & \text{|p_1| or |p_2| not a variable} \\
\rulename{dsparr} &  \ensuremath{\textbf{ar}\{\Varid{p}_{\mathrm{1}};\,\mydots;\,\Varid{p}_{\Varid{n}}\}\mathrel{\coloneqq}\Varid{t}} & \rightarrow &
      \begin{array}[t]{@{}l@{}}
        \ensuremath{\Varid{p}_{\mathrm{1}}\mathrel{\coloneqq}\Varid{x}_{\mathrm{1}};\,\mydots;\,\Varid{p}_{\Varid{n}}\mathrel{\coloneqq}\Varid{x}_{\Varid{n}}} \\
        \ensuremath{\textbf{ar}\{\Varid{x}_{\mathrm{1}}\mathbin{:}\textbf{any};\,\mydots;\,\Varid{x}_{\Varid{n}}\mathbin{:}\textbf{any}\}\hspace{-0.14em}=\hspace{-0.14em}\Varid{t};\,}
      \end{array}  & \text{none of \ensuremath{\Varid{p}_{\Varid{i}}} is \ensuremath{\mathinner{\ldotp\ldotp}\Varid{p}}} \\[1mm]
\rulename{dspamp1} & \ensuremath{ xs \;\!\textbf{\&}\!\;\Varid{x}\mathrel{\coloneqq}\Varid{t}} & \rightarrow & \ensuremath{\textbf{ar}\{ xs \;\!\textbf{\&}\!\;\Varid{x}\mathrel{\coloneqq}\Varid{t}\}} & \text{outside an array}\\
\rulename{dspamp2} & \ensuremath{\textbf{ar}\{\mydots;\, xs \;\!\textbf{\&}\!\;\Varid{x}\mathrel{\coloneqq}\Varid{t};\,\mydots\}} & \rightarrow & \ensuremath{\textbf{ar}\{\mydots;\, xs \mathrel{\coloneqq}\Varid{t};\,\Varid{x}\mathrel{\coloneqq}\Varid{t};\,\mydots\}} & \text{inside an array}\\
\end{array} \\
\\
\begin{array}{lrcll}
\multicolumn{5}{l}{\textbf{Application and unification}} \\
% \rulename{dswhere} &  |t_1 where t_2| & \rightarrow & |x := t_1 ; t_2; x| & \text{fresh |x|}\\
% Don't desugar where
\rulename{dsapp}   & \ensuremath{\Varid{t}_{\mathrm{1}} (\Varid{t}_{\mathrm{2}})} & \rightarrow & \ensuremath{{\bf chk}\langle\textbf{succeeds}\rangle\{\Varid{t}_{\mathrm{1}} [\mskip1.5mu \Varid{t}_{\mathrm{2}}\mskip1.5mu]\}} \\
\\
\multicolumn{5}{l}{\textbf{Array splice}} \\
\rulename{dssplice} & \multicolumn{3}{l}{\quad\ensuremath{\textbf{ar}\{\Varid{t}_{\mathrm{1}};\,\mydots;\,t_{k-1};\,\mathinner{\ldotp\ldotp}\Varid{t}_{\Varid{k}};\,t_{k+1};\,\mydots;\,\Varid{t}_{\Varid{n}}\}}} \\
&  & \rightarrow & \multicolumn{2}{l}{\ensuremath{ \textbf{append}\;[\mskip1.5mu \textbf{ar}\{\Varid{t}_{\mathrm{1}};\,\mydots;\,t_{k-1}\},\textbf{append} [\mskip1.5mu \Varid{t}_{\Varid{k}},\textbf{ar}\{t_{k+1};\,\mydots;\,\Varid{t}_{\Varid{n}}\}\mskip1.5mu]\mskip1.5mu]}}\\
&  & \multicolumn{3}{l}{\text{where \ensuremath{\textbf{append}} is a reversible function appending two arrays}} \\
\\
\multicolumn{5}{l}{\textbf{Array notation}} \\
\rulename{dsunit} &  \ensuremath{()} & \rightarrow & \ensuremath{\textbf{ar}\{\}} \\
\rulename{dstuple} &  \ensuremath{\Varid{t}_{\mathrm{1}},\mydots,\Varid{t}_{\Varid{n}}} & \rightarrow & \ensuremath{\textbf{ar}\{\Varid{t}_{\mathrm{1}};\,\mydots;\,\Varid{t}_{\Varid{n}}\}} \\
\\
\multicolumn{5}{l}{\textbf{Map notation}} \\
\rulename{dsmap}  &  \ensuremath{\textbf{map}\{\Varid{t}_{\mathrm{1}},\mydots,\Varid{t}_{\Varid{n}}\}} & \rightarrow & \multicolumn{2}{l}{\ensuremath{\textbf{mkMap} (\textbf{for}\;(\Varid{f}\mathbin{:}\textbf{ar}\{\Varid{t}_{\mathrm{1}};\,\mydots;\,\Varid{t}_{\Varid{n}}\};\,\Varid{i}\!\rightarrow\!\Varid{a}\mathbin{:}\Varid{f})\;\{\mskip1.5mu (\Varid{i},\Varid{a})\mskip1.5mu\})}} \\
                  &                           &             & & \text{fresh \ensuremath{\Varid{f}},\ensuremath{\Varid{i}},\ensuremath{\Varid{a}}}
\\
\multicolumn{5}{l}{\textbf{Function notation}} \\
\rulename{dstype} &  \ensuremath{\textbf{type} \{\mskip1.5mu \Varid{t}\mskip1.5mu\}} & \rightarrow & \ensuremath{\textbf{fun}_{\textbf{c}}(\Varid{x}\mathrel{\coloneqq}\Varid{t})\langle\textbf{succeeds}\rangle\{\Varid{x}\}} & \text{fresh \ensuremath{\Varid{x}}}\\

\rulename{dsfarrow} &  \ensuremath{\Varid{t}_{\mathrm{1}}\;\!=>\!\!\;\Varid{t}_{\mathrm{2}}} & \rightarrow & \ensuremath{\textbf{fun}_{\textbf{c}}(\Varid{t}_{\mathrm{1}})\langle\textbf{succeeds}\rangle\{\Varid{t}_{\mathrm{2}}\}} \\
\rulename{dsfun1} &  \ensuremath{\textbf{fun}_{\Varid{oc}}(\Varid{a}_{\mathrm{1}}\;\Varid{a}_{\mathrm{2}}\;\mydots\;\Varid{a}_{\Varid{n}})\{\Varid{t}\}} & \rightarrow & \ensuremath{\textbf{fun}_{\Varid{oc}}(\Varid{a}_{\mathrm{1}})\{\textbf{fun}_{\Varid{oc}}(\Varid{a}_{\mathrm{2}}\;\mydots\;\Varid{a}_{\Varid{n}})\{\Varid{t}\}\}} \\
\rulename{dsfun2a}  &   \ensuremath{\textbf{fun}_{\Varid{oc}}(\Varid{a}\langle\omega\rangle)\{\Varid{b}\}} & \rightarrow & \ensuremath{\textbf{fun}_{\Varid{oc}}(\Varid{a})\langle\omega\rangle\{\Varid{b}\}} & \\
\rulename{dsfun2b}  &   \ensuremath{\textbf{fun}_{\Varid{oc}}(\Varid{a})\{\Varid{b}\}} & \rightarrow & \ensuremath{\textbf{fun}_{\Varid{oc}}(\Varid{a})\langle\textbf{succeeds}\rangle\{\Varid{b}\}} & \\
% \rulename{dsfun3}  &   |functionOCX oc a (eff fx) b| & \rightarrow & |functionOC oc a (check fx b)| & \\
\\
\multicolumn{5}{l}{\textbf{Conditional and for-loop notation}} \\
\rulename{dsif1} & \ensuremath{\textbf{if}\;\{\mskip1.5mu \Varid{e}\mskip1.5mu\}} & \rightarrow & \ensuremath{\textbf{if}\;(\Varid{e})\;\textbf{else}\;\textbf{false}} \\
\rulename{dsif2} &  \ensuremath{\textbf{if}\;(\Varid{e}_{\mathrm{1}})\;\textbf{then}\;\Varid{e}_{\mathrm{2}}} & \rightarrow & \ensuremath{\textbf{if}\;(\Varid{e}_{\mathrm{1}})\;\textbf{then}\;\Varid{e}_{\mathrm{2}}\;\textbf{else}\;\textbf{false}} \\
\rulename{dsif3} &  \ensuremath{\textbf{if}\;(\Varid{e}_{\mathrm{1}})\;\textbf{else}\;\Varid{e}_{\mathrm{2}}} & \rightarrow & \ensuremath{\textbf{if}\;(\Varid{x}\mathrel{\coloneqq}\Varid{e}_{\mathrm{1}})\;\textbf{then}\;\Varid{x}\;\textbf{else}\;\Varid{e}_{\mathrm{2}}} & \text{\ensuremath{\Varid{x}} fresh} \\
\rulename{dsfor} &  \ensuremath{\textbf{for}\;\{\mskip1.5mu \Varid{e}\mskip1.5mu\}} & \rightarrow & \ensuremath{\textbf{for}\;(\Varid{x}\mathrel{\coloneqq}\Varid{e})\;\{\mskip1.5mu \Varid{x}\mskip1.5mu\}} & \text{\ensuremath{\Varid{x}} fresh} \\
\end{array}
\end{array}
$$
\caption{Source Verse to Essential Verse rewriting (1)} \label{fig:sv-to-sv-1}
\end{figure}

\begin{figure}
$$
\begin{array}{lrcll}
  \multicolumn{5}{l}{\textbf{Operators}} \\
\rulename{dspreq} &  \ensuremath{\mathbin{?}\Varid{e}}     & \rightarrow & \ensuremath{\textbf{type}\;\{\mskip1.5mu \textbf{false}\hspace{0.19em}\mathop{\rule[-0.1em]{0.15em}{0.75em}}\hspace{0.2em}\mathbf{let}\;(\Varid{x}\mathbin{:}\Varid{e})\;\{\mskip1.5mu \textbf{truth}\{\Varid{x}\}\mskip1.5mu\}\mskip1.5mu\}} & \text{\ensuremath{\Varid{x}} fresh} \\
\rulename{dspostq} & \ensuremath{\Varid{e}\mathbin{?}}     & \rightarrow & \ensuremath{\Varid{e} [\mskip1.5mu \anonymous \mskip1.5mu]} \\
\rulename{dsop1} &   \ensuremath{\Varid{op}_{\Varid{b}}\;\Varid{e}} & \rightarrow & \ensuremath{\textbf{prefix}\textquote \Varid{op}_{\Varid{b}}\textquote\;[\mskip1.5mu \Varid{e}\mskip1.5mu]} & \text{if \ensuremath{\Varid{op}_{\Varid{b}}} can fail}\\
\rulename{dsop2} &   \ensuremath{\Varid{op}_{\Varid{b}}\;\Varid{e}} & \rightarrow & \ensuremath{\textbf{prefix}\textquote \Varid{op}_{\Varid{b}}\textquote\;(\Varid{e})} \\
\rulename{dsop3} &   \ensuremath{\Varid{e}\;\Varid{op}_{\Varid{a}}} & \rightarrow & \ensuremath{\textbf{postfix}\textquote \Varid{op}_{\Varid{a}}\textquote\;[\mskip1.5mu \Varid{e}\mskip1.5mu]} & \text{if \ensuremath{\Varid{op}_{\Varid{a}}} can fail}\\
\rulename{dsop4} &   \ensuremath{\Varid{e}\;\Varid{op}_{\Varid{a}}} & \rightarrow & \ensuremath{\textbf{postfix}\textquote \Varid{op}_{\Varid{a}}\textquote\;(\Varid{e})} \\
\rulename{dsop5} &   \ensuremath{\Varid{e}_{\mathrm{1}}\;\Varid{op}_{\Varid{i}}\;\Varid{e}_{\mathrm{2}}} & \rightarrow & \ensuremath{\textbf{operator}\textquote \Varid{op}_{\Varid{i}}\textquote\;[\mskip1.5mu \Varid{e}_{\mathrm{1}},\Varid{e}_{\mathrm{2}}\mskip1.5mu]} & \text{if \ensuremath{\Varid{op}_{\Varid{i}}} can fail}\\
\rulename{dsop6} &   \ensuremath{\Varid{e}_{\mathrm{1}}\;\Varid{op}_{\Varid{i}}\;\Varid{e}_{\mathrm{2}}} & \rightarrow & \ensuremath{\textbf{operator}\textquote \Varid{op}_{\Varid{i}}\textquote\;(\Varid{e}_{\mathrm{1}},\Varid{e}_{\mathrm{2}})} \\
\rulename{dsand} &   \ensuremath{\Varid{e}_{\mathrm{1}}\;\textbf{and}\;\Varid{e}_{\mathrm{2}}} & \rightarrow & \ensuremath{\Varid{e}_{\mathrm{1}};\,\Varid{e}_{\mathrm{2}}} \\
\rulename{dsor}  &   \ensuremath{\Varid{e}_{\mathrm{1}}\;\textbf{or}\;\Varid{e}_{\mathrm{2}}} & \rightarrow & \ensuremath{\textbf{if}\;(\Varid{e}_{\mathrm{1}})\;\textbf{else}\;\textbf{if}\;(\Varid{e}_{\mathrm{2}})\;\textbf{else}\mathbin{:}\textbf{false}} \\
\rulename{dsbang} &   \ensuremath{\mathbin{!}\Varid{e}} & \rightarrow & \ensuremath{\textbf{if}\;\Varid{e}\;\textbf{then}\mathbin{:}\textbf{false}\;\textbf{else}\;\textbf{false}} \\
  \\
  \multicolumn{5}{l}{\textbf{Let, block, case}} \\
\rulename{dslet}   &    \ensuremath{\mathbf{let}\;(\Varid{e})\;\textbf{in}\;\Varid{b}} & \rightarrow & \ensuremath{\Varid{e};\,\Varid{b}} & \textit{see note} \\
\rulename{dsblock} &    \ensuremath{\textbf{block}\;\Varid{b}} & \rightarrow & \ensuremath{\Varid{b}} & \textit{see note}  \\
\rulename{dscase1} &    \ensuremath{\textbf{case}\;\{\mskip1.5mu \Varid{e}_{\mathrm{1}};\,\mydots\;\Varid{e}_{\Varid{n}}\mskip1.5mu\}} & \rightarrow & \ensuremath{(\Varid{x}\mathbin{:}\textbf{any})\Rightarrow \textbf{case}\;(\Varid{x})\;\{\mskip1.5mu \Varid{e}_{\mathrm{1}};\,\mydots\;\Varid{e}_{\Varid{n}}\mskip1.5mu\}} & \text{\ensuremath{\Varid{x}} fresh} \\
\rulename{dscase2} &    \ensuremath{\textbf{case}\;(\Varid{e})\;\{\mskip1.5mu \Varid{e}_{\mathrm{1}};\,\mydots\;\Varid{e}_{\Varid{n}}\mskip1.5mu\}} & \rightarrow & \ensuremath{\Varid{x}\mathrel{\coloneqq}\Varid{e};\,\Varid{e}_{\mathrm{1}} [\mskip1.5mu \Varid{x}\mskip1.5mu]\;\textbf{or}\;\mydots\;\textbf{or}\;\Varid{e}_{\Varid{n}} [\mskip1.5mu \Varid{x}\mskip1.5mu]} & \text{\ensuremath{\Varid{x}} fresh} \\
  \multicolumn{5}{l}{\lennart{This desugaring does not work for multivalued \ensuremath{\Varid{e}_{\Varid{i}}}.}} \\
  \\
  \multicolumn{5}{l}{\textbf{Misc}} \\
\rulename{dsopt1} &   \ensuremath{\textbf{option} \{\mskip1.5mu \mskip1.5mu\}} & \rightarrow & \ensuremath{\textbf{false}} \\
\rulename{dsopt2} &   \ensuremath{\textbf{option} \{\mskip1.5mu \Varid{e}\mskip1.5mu\}} & \rightarrow & \ensuremath{\textbf{if}\;(\Varid{x}\mathrel{\coloneqq}\Varid{e})\;\textbf{then}\;\textbf{truth}\{\Varid{x}\}\;\textbf{else}\;\textbf{false}} & \text{\ensuremath{\Varid{x}} fresh} \\
% \rulename{dstruth} &  |truth (e}| & \rightarrow & |amap(e eqgt e)| \\
\end{array}
$$
{\em Note}: The desugaring assumes a scope correct Verse program.
This means that there can be no variable name clashes in \ensuremath{\mathbf{let}} and \ensuremath{\textbf{block}}.
pA more robust desugaring could \textalpha-convert such clashes.
\caption{Source Verse to Source Verse rewriting (2)} \label{fig:sv-to-sv-2}
\end{figure}
